\subsection{GPT-5.3-Codex mit Basis Prompt}
In diesem Abschnitt werden die Ergebnisse des Experiments –– GPT-5.3-Codex mit Basis Prompt –– dargestellt (\ref{tab:GPT-5.3-Codex_basis_prompt_evaluation}).
Die Bewertung erfolgt anhand der zuvor definierten Evaluationskriterien.
Alle Validierungs- und Testverfahren werden mithilfe eines separaten Prompts erstellt und durchgeführt. Dadurch können die implementierten Funktionen nach jedem Entwicklungsschritt systematisch überprüft werden. Dieses Vorgehen stellt sicher, dass neu entwickelte Komponenten die definierten Anforderungen erfüllen und dass bestehende Funktionalitäten durch Änderungen nicht unbeabsichtigt beeinträchtigt werden. Auf diese Weise wird eine kontinuierliche Qualitätssicherung während des gesamten Entwicklungsprozesses gewährleistet.
Darüber hinaus wurden zusätzliche Domänenfunktionen wie \textit{Order Item}, \textit{Occupancy Items}, \textit{Services}, Kommentare und Raumstandorte erfolgreich umgesetzt.


\newpage
\begin{table}[H]
\centering
\small
\caption{Ergebnisse von GPT-5.3-Codex mit Basis Prompt}
\label{tab:GPT-5.3-Codex_basis_prompt_evaluation}
\renewcommand{\arraystretch}{1.15}

\begin{tabularx}{\textwidth}{@{}>{\RaggedRight\arraybackslash}p{5.0cm} >{\centering\arraybackslash}p{1.4cm} >{\RaggedRight\arraybackslash}X@{}}
\toprule
\textbf{Evaluationskriterium} & \textbf{Wert} & \textbf{Beobachtung} \\
\midrule
\multicolumn{3}{@{}l@{}}{\textbf{Funktionale Kriterien}} \\
\addlinespace[0.2em]
Allgemeine Funktionalität & 1 &
Grundlegende Benutzeroberfläche ist vorhanden, jedoch sind erweiterte Funktionen wie Filterung, Pagination sowie vollständige CRUD-Operationen unvollständig implementiert. \\
Geschäftsobjekte und CRUD-Funktionalität & 1 &
\textit{Customers, Orders, Rooms, Tasks und Invoices} wurden umgesetzt. Mehrere Module besitzen keine CRUD-Funktionalität. \\
Erweiterte Domänenfunktionen & 2 &
\textit{\textit{Order Items, Occupancy Items, Services, Kommentare und Room Locations}} wurden vollständig implementiert. \\
Benutzer- und Rechteverwaltung & 0 &
Benutzer- und Rollenmodelle sind nicht vorhanden, jedoch fehlt eine vollständige Rechteverwaltung und Authorisierung. \\
\midrule
\multicolumn{3}{@{}l@{}}{\textbf{Technische Kriterien}} \\
\addlinespace[0.2em]
Modell-Validierungen und Geschäftslogik & 1 &
Modell-Validierungen, Geschäftsregeln und technische Qualitätssicherungsmechanismen wurden teilweise umgesetzt. \\
Unit-Tests & 2 &
Es wurden automatisierte Unit-Tests generiert. \\
Codequalität und Wartbarkeit & 1 &
Die technische Struktur enthält teilweise nachweisbare Mechanismen zur Qualitätssicherung oder Testbarkeit. \\
\midrule
\multicolumn{3}{@{}l@{}}{\textbf{Integrationskriterien}} \\
\addlinespace[0.2em]
Frontend-Backend-Integration & 1 &
Grundlegende Interaktionen zwischen Benutzeroberfläche und Backend sind vorhanden. Komplexe Workflows wurden jedoch nicht vollständig umgesetzt. \\
End-to-End-Workflows & 1 &
Teilweise keine vollständigen Benutzerabläufe oder automatisierten Integrationstests vorhanden. \\
\midrule
\multicolumn{3}{@{}l@{}}{\textbf{Gesamtergebnis}} \\
\addlinespace[0.2em]
Erreichte Gesamtpunktzahl & 10 / 18 &
Das generierte System erreicht 10 von 18 möglichen Bewertungspunkten und erfüllt damit etwa 55,6\% der definierten Evaluationskriterien. \\
\bottomrule
\end{tabularx}
\end{table}



Im Bereich der Qualitätssicherung konnte GPT-5.3-Codex die im Prompt definierten Anforderungen weitgehend erfüllen. Der Basis-Prompt forderte explizit die Erstellung von Unit-Tests für relevante Softwarekomponenten, einschließlich Datenmodellen, Geschäftslogik und Views. Darüber hinaus wurde die Implementierung eines Integrationstests verlangt, der den vollständigen Geschäftsprozess von der Anlage eines Kunden über die Erstellung einer Buchung bis hin zur Generierung einer Rechnung abdeckt.

Die Analyse des generierten Systems zeigte, dass für die wesentlichen Komponenten der Anwendung automatisierte Unit-Tests erzeugt wurden. Insbesondere wurden Tests für Datenmodelle, Controller- beziehungsweise View-Komponenten sowie Teile der Geschäftslogik implementiert. Darüber hinaus enthielt das generierte Projekt automatisierte Modell-Validierungen der definierten Geschäftsprozesse.

Zusätzlich wurde ein Integrationstest bereitgestellt, welcher den im Prompt geforderten End-to-End-Ablauf abbildet. Der Test umfasst die Erstellung eines Kunden, die Anlage einer Buchung sowie die Generierung einer Rechnung und kann automatisiert mit \texttt{pytest} ausgeführt werden. Dadurch wurde nachgewiesen, dass die zentralen Systemfunktionen nicht nur implementiert, sondern auch durch automatisierte Tests abgesichert wurden.

Im Vergleich zu den anderen untersuchten Modellen stellt dies einen wesentlichen Qualitätsvorteil dar, da die Generierung automatisierter Testfälle und vollständiger Integrationsszenarien in den übrigen Experimenten nur teilweise oder gar nicht umgesetzt wurde.

Schwächen zeigen sich hingegen im Bereich der Benutzer- und Rechteverwaltung. Obwohl grundlegende Benutzermodelle vorhanden sind, wurde keine vollständige Rollen- und Autorisierungslogik implementiert. Dies deutet darauf hin, dass sicherheitsrelevante Anforderungen nicht vollständig durch den strukturierten Prompt abgedeckt wurden.

Auch die technischen Kriterien wurden nur teilweise erfüllt. Positiv hervorzuheben ist die Implementierung von Modell-Validierungen und Geschäftslogik sowie die automatische Generierung von Unit-Tests. Dagegen konnten keine ausreichenden Maßnahmen zur Sicherstellung der Codequalität und langfristigen Wartbarkeit identifiziert werden. Dies legt nahe, dass der Fokus der Codegenerierung primär auf der funktionalen Umsetzung und weniger auf architektonischen Qualitätsmerkmalen lag.

Im Bereich der Integration wurden lediglich Teilergebnisse erzielt. Zwar sind grundlegende Interaktionen zwischen Frontend und Backend vorhanden, jedoch fehlen vollständig umgesetzte End-to-End-Workflows sowie umfassende Integrationstests. Dadurch wird die praktische Einsetzbarkeit der Anwendung eingeschränkt, da das Zusammenspiel einzelner Komponenten nicht durchgängig abgesichert ist.

Insgesamt erreicht das generierte System 10 von 18 möglichen Bewertungspunkten und erfüllt damit etwa 55,6\% der definierten Evaluationskriterien. Die Ergebnisse verdeutlichen, dass GPT-5.3-Codex mit Basis-Prompts in der Lage ist, umfangreiche funktionale Anforderungen sowie die grundlegende Geschäftslogik erfolgreich umzusetzen. Gleichzeitig bestehen Defizite in den Bereichen Autorisierung, Wartbarkeit, Integration und Performance, wodurch zusätzliche Entwicklungs- und Qualitätssicherungsmaßnahmen erforderlich sind, um ein produktionsreifes System zu erhalten.